% ============================================================
% CAPITOLO 3 — Il Metodo ADLC
% ============================================================
\chapter{Il Metodo ADLC: Come si Lavora in 0-Code}

\section*{Cosa Imparerai}
Alla fine di questo capitolo saprai:
\begin{itemize}
  \item Le 7 fasi dell'Agent Development Life Cycle applicate allo sviluppo di applicazioni
  \item Come scrivere un file \texttt{\_CONTEXT.md} efficace
  \item Cos'è il Context Engineering e perché è la competenza chiave del paradigma 0-code
  \item Come strutturare un progetto per massimizzare la qualità dell'output dell'IA
\end{itemize}

% ---------------------------------------------------------
\section{Il Ciclo di Vita: Le 7 Fasi dell'ADLC}

L'ADLC (Agent Development Life Cycle) è il metodo di lavoro strutturato per sviluppare con l'IA. Non è teoria accademica --- è la sequenza di passi che seguirai in ogni progetto di questo libro.

\begin{small}
\begin{longtable}{c l l l}
\toprule
\textbf{Fase} & \textbf{Nome} & \textbf{Cosa fai} & \textbf{Deliverable} \\
\midrule
\endhead
\textbf{0} & Preparazione & Definisci il problema, valuti adeguatezza & Ipotesi di soluzione \\
\textbf{1} & Inquadramento & Stabilisci confini, tecnologie, vincoli & Requisiti in linguaggio naturale \\
\textbf{2} & Definizione & Scrivi \texttt{\_CONTEXT.md} e \texttt{SKILL.md} & Documenti di contesto \\
\textbf{3} & Simulazione & Testi con un piccolo caso & Proof of Value \\
\textbf{4} & Implementazione & L'IA genera codice, tu supervisioni & Codice funzionante \\
\textbf{5} & Rilascio & Test, fix, deploy & Applicazione in produzione \\
\textbf{6} & Apprendimento & Aggiorni il contesto & Contesto migliorato \\
\bottomrule
\end{longtable}
\end{small}

\begin{notebox}
\textbf{Verso il framework professionale.} In questo capitolo vedremo un contesto base con un singolo \texttt{\_CONTEXT.md}. Se ti stai chiedendo come questo metodo scali su team reali e architetture complesse, nell'\textbf{Appendice~E} analizzeremo un framework ADLC professionale da 17~file usato in produzione --- con moduli per fase, NemoClaw e governance infrastrutturale. I principi che impari qui sono gli stessi; cambiano solo scala e sofisticazione.
\end{notebox}

\subsection{Come si applica nella pratica quotidiana}

Per un progetto semplice (Hello World, CLI tool), le fasi 0--3 si comprimono in 5~minuti di riflessione. Per un progetto complesso (web app full-stack), ogni fase ha un suo peso specifico.

\textbf{Esempio: costruire un'API REST per un blog}
\begin{itemize}
  \item \textbf{Fase~0}: ``Mi serve un backend per il mio blog, con post, commenti e autenticazione''
  \item \textbf{Fase~1}: ``Userò Node.js con Express, PostgreSQL con Prisma, autenticazione JWT''
  \item \textbf{Fase~2}: Scrivo il \texttt{\_CONTEXT.md} con architettura, struttura, vincoli
  \item \textbf{Fase~3}: Chiedo all'IA di generare solo l'endpoint \texttt{/health} per verificare
  \item \textbf{Fase~4}: Procedo con tutti gli endpoint, modello dati, middleware
  \item \textbf{Fase~5}: Test, fix, deploy su Railway
  \item \textbf{Fase~6}: Noto che l'IA ha generato password in plain text $\rightarrow$ aggiorno il contesto
\end{itemize}

\begin{quote}
La Fase~6 è cruciale: il \texttt{\_CONTEXT.md} è un \textbf{documento vivente} che migliora a ogni progetto.
\end{quote}

% ---------------------------------------------------------
\subsection{Quando l'IA Deraglia: Gestire i Fallimenti}

Prima o poi succederà: l'IA entrerà in un \textbf{loop di errori} --- corregge un bug, ne introduce un altro, tenta di risolvere il secondo e rompe il codice che funzionava. È un comportamento noto come ``loop di allucinazioni'' ed è importante sapere come gestirlo.

\begin{warnbox}
\textbf{Le 5 tecniche di recupero quando l'IA si incarta:}
\begin{enumerate}
  \item \textbf{Ferma e resetta la conversazione.} Chiudi la chat corrente e aprine una nuova. Spesso l'IA accumula contesto errato nella memoria della sessione, e una chat pulita risolve il problema.
  \item \textbf{Torna all'ultima versione funzionante.} Se usi Git (e dovresti), esegui \texttt{git stash} o \texttt{git checkout .} per eliminare tutte le modifiche non salvate. Se non usi Git, tieni sempre una copia di backup prima di sessioni complesse.
  \item \textbf{Spezza il task in parti più piccole.} Invece di ``Implementa l'autenticazione completa'', chiedi prima ``Genera solo il middleware di verifica JWT'', verifica che funzioni, poi procedi con il passo successivo.
  \item \textbf{Sii più specifico nel contesto.} Se l'IA genera codice sbagliato, spesso il problema è nel \texttt{\_CONTEXT.md}. Aggiungi un vincolo esplicito: ``NON usare CommonJS. Usa SOLO import/export ESModules.''
  \item \textbf{Cambia modello o strumento.} Se Copilot non riesce, prova Claude Code (o viceversa). Modelli diversi hanno punti di forza diversi.
\end{enumerate}
\end{warnbox}

\begin{tipbox}
\textbf{Regola d'oro}: se l'IA ha fallito due volte sullo stesso problema con lo stesso approccio, \textbf{non insistere}. Cambia strategia: spezza il task, aggiungi contesto, o resetta la conversazione. Insistere produce solo codice sempre più fragile.
\end{tipbox}

\subsection{Fase~6 in Pratica: Come un Vincolo Salva il Progetto}

L'apprendimento continuo (Fase~6) non è un'attività teorica. È il meccanismo che trasforma errori in difese permanenti. Ecco tre scenari reali dove un \texttt{\_CONTEXT.md} lacunoso ha causato problemi --- e come una singola riga di vincolo li ha risolti.

\begin{dangerbox}
\textbf{Disastro~1: Password salvate in chiaro.} L'IA genera un endpoint di registrazione utente e salva la password nel database come testo semplice. L'output ha superato i test funzionali (la registrazione funziona), ma un audit di sicurezza rivelerebbe tutte le password visibili.\\[4pt]
\textbf{Il vincolo che risolve}: aggiungere al \texttt{\_CONTEXT.md}:\\
\texttt{SEC-01: Le password DEVONO essere hashate con bcrypt (costo~12). Mai salvare password in chiaro.}\\[4pt]
Da quel momento, \textbf{ogni} generazione successiva usa bcrypt automaticamente.
\end{dangerbox}

\begin{dangerbox}
\textbf{Disastro~2: L'IA cambia stack a metà progetto.} Stai lavorando con ESModules (\texttt{import/export}) e l'IA, in una nuova sessione di chat, genera codice CommonJS (\texttt{require/module.exports}). Il progetto smette di compilare.\\[4pt]
\textbf{Il vincolo che risolve}:\\
\texttt{Linguaggio: JavaScript ES2022+ con ESModules. NON usare require() o module.exports. Solo import/export.}\\[4pt]
Senza questo vincolo l'IA ``dimentica'' la convenzione ad ogni nuova sessione.
\end{dangerbox}

\begin{dangerbox}
\textbf{Disastro~3: Dipendenze fantasma.} L'IA installa una libreria che non era prevista (es.\ \texttt{lodash}) quando bastavano le funzioni native. Il progetto accumula dipendenze inutili, aumentando superficie d'attacco e dimensione del bundle.\\[4pt]
\textbf{Il vincolo che risolve}:\\
\texttt{Dipendenze: SOLO quelle elencate nel \_CONTEXT.md. Non aggiungere nuove dipendenze senza approvazione esplicita.}
\end{dangerbox}

\begin{tipbox}
\textbf{Dopo ogni errore}, poniti la domanda: ``Quale riga avrei potuto aggiungere al \texttt{\_CONTEXT.md} per prevenire questo problema?'' Poi aggiungila. È l'essenza della Fase~6.
\end{tipbox}

% ---------------------------------------------------------
\section{Context Engineering: L'Arte di Comunicare con l'IA}\index{Context Engineering}

Il Context Engineering è la disciplina fondamentale dello sviluppo 0-code. È la differenza tra un'IA che genera codice mediocre e un'IA che genera software di qualità professionale.

\subsection{Il Problema dell'Amnesia}

I modelli IA (LLM) hanno una proprietà fondamentale: sono \textbf{stateless} --- non hanno memoria tra le sessioni. Ogni volta che apri una nuova finestra di chat o riavvii Copilot, l'IA parte da zero. Non ricorda:
\begin{itemize}
  \item Le decisioni architetturali prese ieri
  \item Le convenzioni di naming del tuo progetto
  \item I bug che avete già risolto insieme
  \item Lo stack tecnologico scelto
\end{itemize}

Se non le fornisci queste informazioni, l'IA \textbf{inventerà} le proprie (un fenomeno chiamato ``allucinazione''). Potrebbe usare MongoDB quando il tuo database è PostgreSQL, o generare codice in CommonJS quando il progetto usa ESModules.

\subsection{La Soluzione: Il file \texttt{\_CONTEXT.md}}\index{CONTEXT.md@\texttt{\_CONTEXT.md}}

Il file \texttt{\_CONTEXT.md} (o \texttt{AGENTS.md}\index{AGENTS.md}) è il \textbf{contratto} tra te e l'IA~\cite{agents-md}. Viene letto automaticamente dall'agente all'inizio di ogni sessione. Contiene tutto ciò che l'IA deve sapere per lavorare correttamente nel tuo progetto.

\begin{quote}
Pensa al \texttt{\_CONTEXT.md} come al brief che daresti a un nuovo sviluppatore il suo primo giorno di lavoro: ``Ecco come è strutturato il progetto, queste sono le tecnologie che usiamo, queste le convenzioni da seguire, e queste le cose da NON fare mai.''
\end{quote}

% ---------------------------------------------------------
\section{Anatomia di un \texttt{\_CONTEXT.md}}

Un buon file di contesto ha 5 sezioni:

\subsection{1.\ Identità e Scopo}

\begin{lstlisting}[language={}]
# Progetto: BlogAPI

Stai lavorando su un backend REST API per un blog personale.
Tecnologie: Node.js 20, Express.js, PostgreSQL 16, Prisma ORM.
Ambiente: sviluppo locale su Windows/macOS.
\end{lstlisting}

Dice all'IA \textbf{cosa} è il progetto e \textbf{con cosa} lavora.

\subsection{2.\ Architettura e Struttura}

\begin{lstlisting}[language={}]
## Struttura del Progetto

src/
  routes/          -> Definizioni degli endpoint (solo routing)
  controllers/     -> Logica dei controller
  services/        -> Business logic (interazione con DB)
  middleware/      -> Auth, error handling, logging
  utils/           -> Funzioni helper
prisma/
  schema.prisma    -> Schema database
tests/
  unit/
  integration/
\end{lstlisting}

Dice all'IA \textbf{dove} mettere le cose. Senza questa sezione, l'IA potrebbe creare strutture diverse ogni volta.

\subsection{3.\ Convenzioni e Standard}

\begin{lstlisting}[language={}]
## Convenzioni

- Naming: camelCase per variabili e funzioni, PascalCase per classi
- Risposte API: sempre { success: boolean, data: T, error?: string }
- Validazione: usa Zod per la validazione degli input
- Errori: gestiti centralmente tramite middleware errorHandler
- Async: usa async/await, mai callback
- Import: usa ESModules (import/export), mai require()
\end{lstlisting}

Dice all'IA \textbf{come} scrivere il codice.

\subsection{4.\ Vincoli e Divieti}

\begin{lstlisting}[language={}]
## Vincoli (OBBLIGATORI)

- NON usare mai query SQL raw. Usa SEMPRE Prisma ORM.
- NON salvare mai password in chiaro. Usa bcrypt con salt >= 12.
- NON esporre mai stack trace negli errori di produzione.
- NON usare `any` in TypeScript. Definisci sempre i tipi.
- NON installare dipendenze senza motivo documentato.
\end{lstlisting}

Dice all'IA cosa \textbf{non deve mai fare}. Questa è la sezione più importante per la sicurezza e la qualità.

\subsection{5.\ Comandi e Workflow}

\begin{lstlisting}[language={}]
## Comandi

- Avviare il server: npm run dev
- Eseguire i test: npm test
- Migrare il database: npx prisma migrate dev
- Generare il client Prisma: npx prisma generate

## Workflow di Sviluppo

Quando implementi una nuova funzionalita':
1. Crea/modifica lo schema Prisma se serve
2. Genera la migrazione
3. Implementa service -> controller -> route
4. Aggiungi test per il nuovo endpoint
5. Verifica che tutti i test passino
\end{lstlisting}

Dice all'IA \textbf{come lavorare} nel progetto.

% ---------------------------------------------------------
\section{Il Tuo Primo \texttt{\_CONTEXT.md} Completo}

Ecco un esempio completo che useremo come base per i prossimi capitoli:

\begin{praticabox}[title={Pratica --- Crea il tuo primo file di contesto}]
Crea un file \texttt{\_CONTEXT.md} nella root di un nuovo progetto:
\end{praticabox}

\begin{lstlisting}[language={},title={\texttt{\_CONTEXT.md} per TaskMaster CLI}]
# Progetto: TaskMaster CLI

## Scopo
Un'applicazione da riga di comando per la gestione di task.
L'utente puo' aggiungere, listare, completare e eliminare task.
I dati vengono salvati in un file JSON locale.

## Tecnologie
- Linguaggio: Python 3.11+
- Librerie: solo moduli standard (argparse, json, os, datetime)
- Nessuna dipendenza esterna ammessa

## Struttura
taskmaster/
+-- _CONTEXT.md
+-- main.py              <- Entry point CLI
+-- task_manager.py      <- Logica di gestione task
+-- storage.py           <- Persistenza su file JSON
+-- tests/
|   +-- test_task_manager.py
|   +-- test_storage.py
+-- data/
    +-- tasks.json       <- Database locale

## Convenzioni
- Naming: snake_case per tutto
- Docstring: ogni funzione pubblica deve avere una docstring
- Type hints: su tutte le firme di funzione
- Output CLI: messaggi chiari e formattati

## Vincoli
- NON usare librerie esterne. Solo standard library.
- NON usare variabili globali.
- NON ignorare gli errori. Gestisci FileNotFoundError, etc.
- Il file tasks.json deve essere creato automaticamente.

## Comandi CLI
- python main.py add "Comprare il latte"
- python main.py list
- python main.py list --status done
- python main.py done 3
- python main.py delete 3

## Testing
- Framework: pytest
- python -m pytest tests/ -v
- Copertura test >= 80%
\end{lstlisting}

\begin{checkpointbox}
Questo file è tutto ciò che serve all'IA per generare un'applicazione CLI completa, testata e funzionante. Nel Capitolo~5 lo metterai in pratica.
\end{checkpointbox}

% ---------------------------------------------------------
\section{Le 7 Regole d'Oro del Context Engineering}\index{Context Engineering!regole}

\subsection*{Regola 1: Sii imperativo, non descrittivo}
\begin{lstlisting}[language={}]
NO:  "Sarebbe preferibile usare async/await"
SI:  "Usa SEMPRE async/await. Mai callback o .then()."
\end{lstlisting}

\subsection*{Regola 2: I divieti sono più importanti dei permessi}
L'IA sa fare mille cose. Il tuo lavoro è dirle le 10 che \textsc{non} deve fare. Ogni divieto nel contesto previene una classe intera di errori.

\subsection*{Regola 3: Mostra la struttura, non descriverla}
\begin{lstlisting}[language={}]
NO: "Organizza i file in modo logico con cartelle separate"
SI: "src/
       controllers/
       models/
       services/"
\end{lstlisting}

\subsection*{Regola 4: Specifica i comandi esatti}
L'IA non deve inventare come avviare il progetto. Se il comando è \texttt{npm run dev}, scrivilo.

\subsection*{Regola 5: Un \texttt{\_CONTEXT.md} per progetto}
Non riusare lo stesso contesto per progetti diversi. Ogni progetto ha il suo.

\subsection*{Regola 6: Aggiorna dopo ogni errore}
Quando l'IA commette un errore che poi correggi, aggiungi un vincolo al \texttt{\_CONTEXT.md} per prevenirlo in futuro. Il contesto è un documento vivente.

\subsection*{Regola 7: Meno è meglio, purché sia preciso}
Un contesto di 50 righe iper-specifiche batte un contesto di 500 righe generiche. L'IA perde efficacia quando il contesto è troppo lungo e dispersivo.

% ---------------------------------------------------------
\section{MCP e Agent Skills: Anteprima}

\subsection{Model Context Protocol (MCP)}\index{MCP}\index{Model Context Protocol|see{MCP}}
Quando l'IA deve interagire con servizi esterni (database, API, Slack, GitHub), utilizza il \textbf{MCP} --- un protocollo standard che collega l'IA agli strumenti~\cite{anthropic:mcp,mcp:spec}. Lo introdurremo nel Capitolo~7 quando collegheremo il backend a PostgreSQL.

\subsection{Agent Skills (\texttt{SKILL.md})}\index{SKILL.md}\index{Agent Skills}
Quando un progetto diventa complesso e ha bisogno di competenze specializzate, puoi creare file \texttt{SKILL.md} che conferiscono all'IA expertise su richiesta. Ne esistono per ogni fase ADLC: dall'\textbf{analisi} (requisiti, user story) al \textbf{design} (valutazione stack, ADR), dall'\textbf{implementazione} (React, Flutter, API, database, sicurezza) fino alle \textbf{operations} (incidenti, monitoring). Li introdurremo nel Capitolo~9 per il frontend React e nel Capitolo~11 per Flutter; l'Appendice~E analizza la copertura completa su tutte le fasi.

Per ora, il \texttt{\_CONTEXT.md} è tutto ciò che ti serve.

% ---------------------------------------------------------
\section{Gestione di Progetti Lunghi: \texttt{PROGRESS.md}}\index{PROGRESS.md}

Per progetti che durano più sessioni, avrai bisogno di un sistema di \textbf{memoria persistente}. L'IA dimentica tutto tra una sessione e l'altra, ma può leggere un file.

Il file \texttt{PROGRESS.md} funge da \textbf{diario del progetto}. Al termine di ogni sessione di lavoro, chiedi all'IA:

\begin{lstlisting}[language={}]
Aggiorna PROGRESS.md con:
1. Cosa e' stato implementato in questa sessione
2. Quali decisioni architetturali sono state prese
3. Quali problemi sono stati incontrati e come risolti
4. Cosa resta da fare
\end{lstlisting}

La prossima sessione, l'IA leggerà questo file e ripartirà con tutta la conoscenza accumulata.

\begin{notebox}
Se usi \textbf{Claude Code}, puoi chiamare questo file \texttt{claude-progress.txt} --- viene letto automaticamente all'inizio di ogni sessione. Con altri strumenti (Copilot, Cursor, Windsurf), istruisci l'agente a leggerlo: \emph{``Leggi PROGRESS.md per capire lo stato del progetto, poi...''}.
\end{notebox}

Approfondiremo questa tecnica nel Capitolo~10.

% ---------------------------------------------------------
\section{Il Workflow Professionale: Dall'Idea all'Implementazione}

Finora hai imparato a scrivere il \texttt{\_CONTEXT.md} manualmente e a chiedere all'IA di implementare. Ma nei progetti reali --- quelli con decine di funzionalità, vincoli di sicurezza e team che si alternano --- serve un passaggio intermedio: \textbf{usare l'IA stessa per decomporre l'idea in unità di lavoro implementabili}.

Il workflow professionale ADLC completo ha sei fasi di decomposizione prima di scrivere una sola riga di codice:

\begin{figure}[H]
\centering
\includegraphics[width=0.5\textwidth]{11-workflow-adlc.png}
\caption{Workflow professionale ADLC: dall'idea all'implementazione in 6 fasi}
\label{fig:workflow-adlc}
\end{figure}

\subsection{Fase 1: Dall'idea alle specifiche}

Parti con un prompt di alto livello. L'IA genera un documento strutturato di requisiti:

\begin{lstlisting}[language={}]
Prompt:
"Devo costruire un'applicazione BookShelf per gestire una libreria
 personale. Backend Node.js + Express, frontend React, mobile Flutter,
 auth OAuth Google.
 Genera un documento di specifiche con requisiti funzionali e non
 funzionali, vincoli tecnici e criteri di accettazione."
\end{lstlisting}

L'IA produce un PRD (Product Requirements Document) con sezioni chiare: obiettivo, utenti target, requisiti funzionali numerati (RF-01, RF-02\ldots), requisiti non funzionali (RNF-01\ldots), vincoli tecnici e priorità.

\begin{tipbox}
Salva il PRD in \texttt{docs/PRD.md} nel progetto. È il documento di riferimento per tutto il resto.
\end{tipbox}

\subsection{Fase 2: Dai requisiti ai casi d'uso}

Prendi i requisiti e chiedi all'IA di derivare i casi d'uso:

\begin{lstlisting}[language={}]
Prompt:
"A partire dal PRD in docs/PRD.md, genera i casi d'uso principali.
 Per ogni caso d'uso indica: attore, precondizione, flusso principale,
 flussi alternativi, postcondizione."
\end{lstlisting}

L'IA produce casi d'uso come:

\begin{lstlisting}[language={}]
## UC-01: Aggiunta libro al catalogo
- Attore: Utente autenticato
- Precondizione: L'utente ha effettuato il login
- Flusso principale:
  1. L'utente clicca "Aggiungi libro"
  2. Compila titolo, autore, anno, genere
  3. Il sistema valida i dati
  4. Il sistema salva il libro nel database
  5. Il sistema mostra il libro nella lista aggiornata
- Flusso alternativo: Se il titolo e' vuoto -> errore di validazione
- Postcondizione: Il libro e' visibile nella libreria dell'utente
\end{lstlisting}

Salva in \texttt{docs/USE\_CASES.md}. Questo documento colma il gap tra ``cosa fa il sistema'' e ``come lo usa una persona reale''.

\subsection{Fase 3: Progettazione e scelte tecnologiche}\index{ADR}\index{Architecture Decision Record|see{ADR}}\index{progettazione}

Questa è la fase in cui \textbf{l'IA propone, ma l'umano decide}. Prima di scomporre il lavoro in epic e task, servono decisioni architetturali che condizioneranno tutto il resto.

\begin{lstlisting}[language={}]
Prompt:
"A partire dal PRD e dai casi d'uso, proponi l'architettura del sistema.
 Per ogni decisione tecnica (database, ORM, autenticazione, struttura
 progetto, pattern architetturali) presenta 2-3 opzioni con pro, contro
 e raccomandazione. Io scelgo, tu documenti in docs/DESIGN.md."
\end{lstlisting}

L'IA genera un documento di design con decisioni strutturate:

\begin{lstlisting}[language={}]
## ADR-01: Scelta ORM

### Contesto
Il backend richiede un ORM per PostgreSQL compatibile con Node.js.
\end{lstlisting}

\begin{center}
\begin{tabularx}{\textwidth}{l X X}
\toprule
\textbf{Opzione} & \textbf{Pro} & \textbf{Contro} \\
\midrule
\textbf{Prisma} & Type-safe, migrazioni auto, ottima DX & Meno flessibile per query complesse \\
\textbf{Drizzle} & Leggero, SQL-like, performante & Community più piccola, meno tutorial \\
\textbf{TypeORM} & Maturo, pattern Active Record & API verbose, manutenzione rallentata \\
\bottomrule
\end{tabularx}
\end{center}

\begin{lstlisting}[language={}]
### Raccomandazione
Prisma -- migliore per un progetto didattico con focus su produttivita'.

### Decisione
-> **[L'UMANO SCEGLIE QUI]**
\end{lstlisting}

L'IA prosegue con le decisioni successive: struttura delle cartelle, strategia di autenticazione, gestione dello stato frontend, pattern API (REST vs.\ GraphQL), strategia di deploy.

\begin{warnbox}
\textbf{Regola fondamentale}: L'IA non deve scegliere da sola le tecnologie. Presenta le opzioni con dati concreti (performance, maturità, curva di apprendimento, costo), ma la scelta finale è sempre dell'umano. Questo vale per:
\begin{itemize}
  \item Stack tecnologico (linguaggio, framework, database)
  \item Pattern architetturali (monolite vs.\ microservizi, REST vs.\ GraphQL)
  \item Servizi cloud (hosting, auth provider, CDN)
  \item Librerie con impatto architetturale (ORM, state management, UI framework)
\end{itemize}
\end{warnbox}

Il risultato è un \texttt{docs/DESIGN.md} con tutte le ADR (Architecture Decision Records) --- decisioni motivate, tracciabili, che l'IA rispetterà in tutte le fasi successive.

Salva in \texttt{docs/DESIGN.md}. Questo documento alimenta il \texttt{\_CONTEXT.md} con le scelte tecnologiche definitive.

\subsection{Fase 4: Dai casi d'uso alle epic}

Le epic raggruppano i casi d'uso in blocchi di lavoro gestibili:

\begin{lstlisting}[language={}]
Prompt:
"A partire dai casi d'uso in docs/USE_CASES.md, organizza il lavoro
 in epic. Ogni epic deve avere: codice (E-001, E-002...), titolo,
 obiettivo, casi d'uso coperti, dipendenze tra epic,
 e ordine di implementazione."
\end{lstlisting}

L'IA produce un indice come:

\begin{center}
\begin{tabularx}{\textwidth}{l l c l}
\toprule
\textbf{Codice} & \textbf{Titolo} & \textbf{Priorità} & \textbf{Dipendenze} \\
\midrule
E-001 & Setup e Infrastruttura & 1 & --- \\
E-002 & Autenticazione OAuth & 2 & E-001 \\
E-003 & CRUD Libri & 3 & E-002 \\
E-004 & Ricerca e Filtri & 4 & E-003 \\
E-005 & Statistiche & 5 & E-003 \\
E-006 & Frontend Web & 6 & E-003 \\
E-007 & App Mobile & 7 & E-003 \\
E-008 & Testing e Sicurezza & 8 & E-006 \\
E-009 & Deploy Produzione & 9 & E-008 \\
\bottomrule
\end{tabularx}
\end{center}

Salva in \texttt{docs/epics/INDEX.md} con un file per epic: \texttt{E-001\_Setup.md}, ecc.

\subsection{Fase 5: Dalle epic ai task}

Ogni epic viene scomposta in task atomici --- l'unità di lavoro che l'IA implementa in una singola sessione:

\begin{lstlisting}[language={}]
Prompt:
"Scomponi l'epic E-003 (CRUD Libri) in task atomici.
 Ogni task deve avere: codice (T-003.1, T-003.2...), titolo,
 acceptance criteria verificabili, file coinvolti,
 vincoli SEC/PERF applicabili, stima in story point (1-5)."
\end{lstlisting}

L'IA genera i task con questo formato:

\begin{lstlisting}[language={}]
# T-003.1 -- Modello Prisma Book
Epic: E-003 | SP: 2 | Stato: TODO | Rischio: LOW

## Obiettivo
Definire il modello Book nello schema Prisma con tutti i campi da PRD.

## Acceptance Criteria
- [ ] Il modello Book e' definito in schema.prisma
- [ ] Include i campi: title, author, year, genre, isbn?, coverUrl?,
      status (enum), rating?, notes?
- [ ] La relazione con User e' configurata (userId, foreign key)
- [ ] La migrazione e' eseguita senza errori
- [ ] Un test verifica la creazione di un record Book

## File coinvolti
- backend/prisma/schema.prisma
- backend/tests/models/book.test.js

## SEC/PERF applicabili
- SEC-02: Ogni query filtra per userId
- PERF-02: Indice su Book.userId e Book.status
\end{lstlisting}

Salva i task in \texttt{docs/epics/tasks/E-003/T-003.1\_Modello\_Book.md}.

\subsection{Fase 6: Implementazione guidata dai task}

Con i task pronti, l'implementazione diventa sistematica:

\begin{lstlisting}[language={}]
Prompt: "Implementa il task T-003.1 (Modello Prisma Book).
 Leggi i dettagli in docs/epics/tasks/E-003/T-003.1_Modello_Book.md.
 Rispetta i vincoli SEC-02 e PERF-02 dal _CONTEXT.md."
\end{lstlisting}

L'IA legge il file del task, conosce gli acceptance criteria, implementa, testa, e segna il task come completato. Poi passi al successivo: \texttt{"Implementa T-003.2."} Ogni sessione chiude uno o più task. Al termine l'IA aggiorna lo stato del task, produce un checkpoint e aggiorna \texttt{\_CONTEXT.md} e \texttt{PROGRESS.md}.

\subsection{Perché questo workflow funziona}

\begin{center}
\begin{tabularx}{\textwidth}{X X}
\toprule
\textbf{Senza decomposizione} & \textbf{Con decomposizione} \\
\midrule
``Costruisci l'app BookShelf'' & ``Implementa T-003.1'' \\
L'IA inventa l'architettura & L'architettura è nelle specifiche \\
Se deraglia, riparti da zero & Se deraglia, riprendi dal task \\
Nessuna tracciabilità & Ogni decisione è documentata \\
Impossibile stimare i tempi & Story point per task \\
\bottomrule
\end{tabularx}
\end{center}

\begin{checkpointbox}
Il workflow Idea $\rightarrow$ Specifiche $\rightarrow$ Use Case $\rightarrow$ Design $\rightarrow$ Epic $\rightarrow$ Task $\rightarrow$ Implementazione è la differenza tra ``sperimentare con l'IA'' e ``fare ingegneria del software con l'IA''. Lo applicherai nel progetto autonomo dell'Appendice~F.
\end{checkpointbox}

% ---------------------------------------------------------
\section*{Riepilogo}

\begin{center}
\begin{tabularx}{\textwidth}{l X}
\toprule
\textbf{Concetto} & \textbf{Cosa hai imparato} \\
\midrule
ADLC & Le 7 fasi del ciclo di vita 0-code \\
\texttt{\_CONTEXT.md} & Come scrivere il contratto con l'IA \\
Context Engineering & Le 7 regole per contesti efficaci \\
Struttura progetto & Le 5 sezioni di un contesto ben scritto \\
MCP / Skills & Anteprima di strumenti avanzati \\
\texttt{PROGRESS.md} & Memoria persistente tra sessioni \\
Workflow professionale & Dall'idea ai task implementabili in 6 fasi \\
\bottomrule
\end{tabularx}
\end{center}
